\subsection{Motivational example}
\subsubsection{Velocity-tracking problem}
    \begin{equation*}
    \left(1 + \cos^2(q)\right)\ddot{q} + g\sin(q) + b\left(\dot{q} + \arctan (\dot{q})\right) = u
    \end{equation*}
    where $q \in \R$ and $\dot{q} \in \R$ are state variables and $u \in \R$ is the torque control input

    By defining $e_1 = \dot{q} - \dot{q}_d$, we get the following error dynamics
    \begin{equation}
        \dot{e}_1 = \gamma(q)u - \gamma(q)g\sin(q) - \gamma(q)b(e_1 + \dot{q}_d) + \arctan(e_1 + \dot{q}) - \ddot{q}
    \end{equation}
    where $\gamma(q) = 1/(1 + \cos^2(q))$. Here, the time-varying function $\gamma(q(t))$ is an uncertain time-varying control coefficient. Let us remark that $\gamma(q(t))$ is bounded by $0 < 1/2 \leq \gamma(q(t)) \leq 1$, with a bounded time derivative.

    Use standard STA
    \begin{equation*}
        \begin{split}
            u_{sta} &= -k_1 \lfloor e_1 \rceil^{1/2} + z \\
            \dot{z} &= -k_2 \lfloor e_1 \rceil^{0}
        \end{split}
    \end{equation*}
    where $\lfloor \cdot \rceil^a = \lvert \cdot \rvert^a sgn(\cdot)$. The closed-loop dynamics then becomes 
    \begin{equation}
        \dot{e}_1 = -k_1 \gamma(q) \lfloor e_1 \rceil^{1/2} + \gamma(q)\left[z - g\right]
    \end{equation}


\subsubsection{Solution}
    \begin{equation}
        \begin{split}
        \dot{e}_2 = &-k_2 \lfloor e_1 \rceil^0 + \dot{\varphi} \underbrace{- bc\gamma k_1 \lfloor e_1 \rceil^{1/2} - bc\gamma k_2 \int_{0}^{t} \lfloor e_1(\tau) \rceil^0 d\tau}_{u(t)} \\
        &+ bc\gamma \int_{0}^{t} \left[\dot{\varphi}(\tau) - b\left(k_1 \lfloor e_1(\tau) \rceil^{1/2}- e_2(\tau)\right)\right]d\tau
        \end{split} 
    \end{equation}
% \begin{frame}[t]
% \frametitle{Motivation}
%     \begin{block}{Problems to be addressed}
%        \begin{enumerate}
%          \item<1-> $k_2$ is designed depending on $k_1$. In some instances this causes and algebraic loop.
%          \item<2-> Although the perturbations $\dot{\varphi}$ can be bounded by a constant, the terms $- bc\gamma k_1 \lfloor e_1 \rceil^{1/2} - bc\gamma k_2 \int_{0}^{t} \lfloor e_1(\tau) \rceil^0 d\tau$ can be very large depending on initial conditions $e_1(0)$ and perturbations.
%          \item<3-> The time variation in the input gain 
%          \item<4->  
%        \end{enumerate}
%     \end{block}
% \end{frame}

\subsection{Generalized super-twisting algorithm}
\subsubsection{Control law}
The generalized super-twisting algorithm \cite{Cas:18} is given by
\begin{equation}
    \begin{split}
        u &= -k_1 \phi_1(\sigma) + z \\ 
        \dot{z} &= -k_2\phi_2(\sigma) \\
        \phi_1(\sigma) &= \lceil \sigma \rfloor^{1/2} + \beta \sigma \\
        \phi_2(\sigma) &= \frac{1}{2} \lceil \sigma \rfloor^0 + \frac{3}{2}\beta\lceil \sigma \rfloor^{1/2} + \beta^2\sigma  
    \end{split}
\end{equation}

\subsection{Finding conditions on gains}
\textbf{Cases}
\begin{itemize}
\item General case
\item State-dependent perturbations and know input gain
\item Time-dependent unknown input gain and perturbation
\item Closing the loop
\end{itemize}

\subsubsection{General case}
Scalar system
\begin{equation}
    \dot{\sigma} = \gamma(x, t)u + \varphi(x, t)
    \label{eq:castillo-general-case}
\end{equation}
\begin{assumption}
    The functions $\gamma(\sigma, t)$ and $\varphi(\sigma, t)$ are Lipschitz continuous functions with respect to $t$, and $\gamma(\sigma, t)$, $\varphi(\sigma, t) \in \mathcal{C}^1$ with respect to $\sigma$.
\end{assumption}
\begin{assumption}
    The uncertain input gain function is bounded by
            \begin{equation}
               0 < k_m \leq \gamma(x, t) \leq k_M
               \label{eq:castillo-7}
            \end{equation}
\end{assumption}
\begin{assumption}
    The perturbation term $\varphi(\sigma, t)$ can be split into two components
         \begin{equation}
            \varphi(\sigma, t) = \varphi_1(\sigma, t) + \varphi_2(\sigma, t),
         \end{equation}
         where is vanishing, i.e. $\varphi_1(0, t) = 0, \ \forall t \leq 0$, and is bounded by
            \begin{equation}
               |\varphi_1(\sigma, t)| \leq \alpha|\phi_1(\alpha)|
            \end{equation}
\end{assumption}
\begin{assumption}
    The total time derivative of $\frac{\varphi_2(\sigma, t)}{\gamma(\sigma, t)}$ is 
            \begin{equation}
               \frac{d}{dt}\left(\frac{\varphi_2(\sigma, t)}{\gamma(\sigma, t)}\right) = 
               \underbrace{\frac{1}{\gamma}\frac{\partial \varphi_2}{\partial t} - \frac{\varphi_2}{\gamma^2}\frac{\partial \gamma}{\partial t}}_{\delta_1(\sigma, t)} + 
               \underbrace{\left(\frac{1}{\gamma}\frac{\partial \varphi_2}{\partial \sigma} - \frac{\varphi_2}{\gamma^2}\frac{\partial \gamma}{\partial \sigma}\right)}_{\delta_2(\sigma, t)}\dot{\sigma}
            \end{equation}
            where $\delta_1(\sigma, t)$ and $\delta_2(\sigma, t)$ are bounded by positive constants
            \begin{equation}
               |\delta_1(\sigma, t)| \leq \bar{\delta}_1, \ |\delta_2(\sigma, t)| \leq \bar{\delta}_2
            \label{eq:castillo-13}
            \end{equation}
\end{assumption}

Closed-loop system
      \begin{equation}
         \dot{\sigma} = -k_1\gamma(\sigma, t)\phi_1(\sigma) + \varphi_1(\sigma, t) + \gamma(\sigma, t)\left(z + \frac{\varphi_2(\sigma, t)}{\gamma(\sigma, t)}\right)
      \end{equation}
      the system can be broken down
      \begin{equation}
         \begin{split}
            \dot{\sigma}_1 &= \gamma(\sigma, t)\left(-k_1\phi_1(\sigma_1) + \frac{\varphi_1(\sigma_1,t)}{\gamma(\sigma, t)} + \sigma_2\right)\\
            \dot{\sigma}_2 &= -k_2\phi_2(\sigma_1)+ \frac{d}{dt}\left(\frac{\varphi_2(\sigma, t)}{\gamma(\sigma, t)}\right)
         \end{split}
      \end{equation}
      where $\sigma_1 = \sigma$ and $\sigma_2 = z + \frac{\varphi_2(\sigma, t)}{\gamma(\sigma, t)}$

   \begin{theorem}[Castillo et al. \cite{Cas:18}]
      \label{th:cas}
      Suppose that $\gamma(\sigma, t)$ and $\varphi(x, t)$ of system \eqref{eq:castillo-general-case} satisfy \eqref{eq:castillo-7} and \eqref{eq:castillo-13}. Then, the states $\sigma_1$ and $\sigma_2$ converge to zero and z converges to $-\varphi(x, t)$, globally and in finite time, if GSTA gains $k_1$, $k_2$ and $\beta > 0$ are design as follows for any $\varepsilon > 0$
      \begin{equation}
         k_1 > \frac{(1+\varepsilon)}{4\varepsilon k_m}\left(\frac{\alpha}{k_m}(k_M - k_m) + \frac{\bar{\delta}_2\varepsilon}{\beta}(k_M + k_m + 2)\right) + \frac{(1+\varepsilon)}{4\varepsilon k_m}\sqrt{\Lambda_{k_1}} + \frac{\alpha}{k_m}
      \end{equation}
      \begin{equation}
         \Lambda_{k_1} = \left(\frac{\alpha}{k_m}(k_M - k_m) + \frac{\bar{\delta}_2\varepsilon}{\beta}(k_M + k_m + 2)\right)^2 + 8\varepsilon(k_M - km) \left(\frac{\bar{\delta}_2}{\beta} + \bar{\delta}_1\right)
      \end{equation}
      \begin{equation}
         k_M \left(2\sqrt{\underbar{h}\epsilon}\bar{c}- \frac{\bar{c}}{2}\sqrt{\Delta_{k_2}}\right)^2 + 2\bar{\delta}_1 < k_2 < k_M \left(2\sqrt{\underbar{h}\epsilon}\bar{c} + \frac{\bar{c}}{2}\sqrt{\Delta_{k_2}}\right)^2 + 2\bar{\delta}_1 
      \end{equation}
      where
      \begin{equation*}
         \Delta_{k_1} =16\underbar{h}\epsilon - 8 \left(\frac{(1+\epsilon)(k_M - k_m)}{\underbar{k}_1 k_m}\right) \times \left(\frac{(1+\epsilon)}{\underbar{k}_1k_m}\left[\frac{\bar{\delta}_2\alpha}{\beta}+2\bar{\delta}_1\right]+\frac{\bar{\delta}_2\epsilon}{\beta} + \frac{\alpha}{k_M}\right),
      \end{equation*}
      \begin{equation*}
         \underbar{h} = 1 - \frac{\bar{\delta}_2}{\beta} \frac{k_m + \epsilon}{(k_1k_m - \alpha)}, \ \bar{c} = \frac{\underbar{k}_1k_m}{(1+\epsilon)(k_M - k_m)}, \ \underbar{k}_1 = k_1 - \frac{\alpha}{k_m}
      \end{equation*}
      Moreover, a system trajectory starting at $\sigma(0) = (\sigma_1(0), |\sigma_2(0))$ reaches the origin in a time smaller than 
      \begin{equation}
         T(\sigma(0)) = \frac{2}{\mu_2} \ln \left(\frac{\mu_2}{\mu_1}V^{1/2}(x_1(0), x_2(0)) + 1\right)
      \end{equation}
   \end{theorem}

Choosing gains
%    \begin{figure}[h]
%       \includegraphics[width=0.6\textwidth]{figures/castillo-fig-1.png}
%    \end{figure}
%    \begin{figure}[h]
%       \includegraphics[width=0.6\textwidth]{figures/castillo-fig-2.png}
%    \end{figure}

\subsection{State-dependent perturbations and know input gain}
System
\begin{equation}
    \dot{\sigma} = u + \varphi(x, t)
    \label{eq:castillo-17}
\end{equation}
\begin{assumption}
    The function $\varphi(\sigma, t)$ is Lipschitz continuous functions with respect to $t$, and $\varphi(\sigma, t) \in \mathcal{C}^1$ with respect to $\sigma$.
\end{assumption}
\begin{assumption}
    The total time derivative of $\frac{\varphi_2(\sigma, t)}{\gamma(\sigma, t)}$ is 
            \begin{equation}
               \frac{d}{dt}\left(\varphi_2(\sigma, t)\right) = 
               \underbrace{\frac{\partial \varphi_2}{\partial t}}_{\delta_1(\sigma, t)} + 
               \underbrace{\left(\frac{\partial \varphi_2}{\partial \sigma}\right)}_{\delta_2(\sigma, t)}\dot{\sigma}
            \end{equation}
\end{assumption}

Closed-loop system
\begin{equation}
    \dot{\sigma} = -k_1\phi_1(\sigma) + \varphi_1(\sigma, t) + \left(z + \frac{\varphi_2(\sigma, t)}{\gamma(\sigma, t)}\right)
\end{equation}
the system can be broken down
\begin{equation}
    \begin{split}
    \dot{\sigma}_1 &= \left(-k_1\phi_1(\sigma_1) + \varphi_1(\sigma_1,t) + \sigma_2\right)\\
    \dot{\sigma}_2 &= -k_2\phi_2(\sigma_1)+ \frac{d}{dt}\left(\varphi_2(\sigma, t)\right)
    \end{split}
\end{equation}
where $\sigma_1 = \sigma$ and $\sigma_2 = z + \varphi_2(\sigma, t)$

\begin{corollary}[1, Castillo et al. \cite{Cas:18}]
    Suppose that $\varphi(\sigma, t)$ of system \eqref{eq:castillo-17} $\delta_1(\sigma, t)$ and $\delta_2(\sigma, t)$ are bounded by $\bar{\delta}_1$ and $\bar{\delta}_2$, respectively. Then, the states $\sigma_1$ and $\sigma_2$ converge to zero and z converges to $-\varphi(\sigma, t)$, globally in finite-time, if GSTA gains $k_1$, $k_2$ and $\beta > 0$ are designed as follows for any $\epsilon > 0$
    \begin{equation}
        k_1 > \frac{2(1+\epsilon)\bar{\delta}_2}{\beta} + \alpha
    \end{equation}
    \begin{equation}
        k_2 > \frac{1}{4\underbar{h}\epsilon}
        \left(
        \frac{(1+\epsilon)}{(k_1 - \alpha)}
        \left(
            \frac{\bar{\delta}_2\alpha}{\beta} + 2\bar{\delta}_1
        \right)
        + \frac{\bar{\delta}_2\epsilon}{\beta} + \alpha
        \right)^2
    \end{equation}
\end{corollary}

\subsection{Time-dependent unknown input gain and perturbation}
Scalar system
\begin{equation}
    \dot{\sigma} = \gamma(t)u + \varphi(t)
    \label{eq:castillo-24}
\end{equation}
\begin{assumption}
    The functions $\gamma(t)$ and $\varphi(t)$ are Lipschitz continuous functions with respect to $t$.
\end{assumption}
\begin{assumption}
    The uncertain input gain function is bounded by
    \begin{equation}
        0 < k_m \leq \gamma(t) \leq k_M
        \label{eq:castillo-25}
    \end{equation}
\end{assumption}
\begin{assumption}
The total time derivative of $\frac{\varphi_2(t)}{\gamma(t)}$ is 
\begin{equation}
    \delta_1(t) := \frac{d}{dt}\left(\frac{\varphi_2(t)}{\gamma(t)}\right) = 
    \frac{1}{\gamma}\frac{\partial \varphi_2}{\partial t} - \frac{\varphi_2}{\gamma^2}\frac{\partial \gamma}{\partial t}
\end{equation}
where $\delta_1(t)$ is bounded by positive a constant
\begin{equation}
    |\delta_1(t)| \leq \bar{\delta}_1
\label{eq:castillo-27}
\end{equation}
\end{assumption}

\subsection{Time-dependent unknown input gain and perturbation}
Closed-loop system
\begin{equation}
    \dot{\sigma} = -k_1\gamma(t)\phi_1(\sigma) + \varphi_1(t) + \gamma(t)\left(z + \frac{\varphi_2(t)}{\gamma(t)}\right)
\end{equation}
the system can be broken down
\begin{equation}
    \begin{split}
    \dot{\sigma}_1 &= \gamma(t)\left(-k_1\phi_1(\sigma_1) + \frac{\varphi_1(t)}{\gamma(t)} + \sigma_2\right)\\
    \dot{\sigma}_2 &= -k_2\phi_2(\sigma_1)+ \frac{d}{dt}\left(\frac{\varphi_2(t)}{\gamma(t)}\right)
    \end{split}
\end{equation}
where $\sigma_1 = \sigma$ and $\sigma_2 = z + \frac{\varphi_2(t)}{\gamma(t)}$

\begin{corollary}[2, Castillo et al. \cite{Cas:18}]
Suppose that $\gamma(t)$ and $\varphi(t)$ of system \eqref{eq:castillo-24} satisfy \eqref{eq:castillo-25} and \eqref{eq:castillo-27}. Then, the states $\sigma_1$ and $\sigma_2$ converge to zero and z converges to $-\frac{\varphi(t)}{\gamma(t)}$, globally and in finite-time, if GSTA gains $k_1, k_2$ and $\beta \geq 0$ are designed as follows for all $\epsilon > 0$
\begin{equation}
    k_1 > \frac{(1+\varepsilon)}{4\varepsilon k_m}\left(\frac{\alpha}{k_m}(k_M - k_m)\right) + \frac{(1+\varepsilon)}{4\varepsilon k_m}\sqrt{\Lambda_{k_1}} + \frac{\alpha}{k_m}
\end{equation}
\begin{equation}
    \Lambda_{k_1} = 
    \left(
    \frac{\alpha}{k_m}(k_M - k_m)
    \right)^2 + 16\varepsilon(k_M - km)\bar{\delta}_1
\end{equation}
\begin{equation}
    k_M \left(2\sqrt{\epsilon}\bar{c}- \frac{\bar{c}}{2}\sqrt{\Delta_{k_2}}\right)^2 + 2\bar{\delta}_1 < k_2 < k_M \left(2\sqrt{\epsilon}\bar{c} + \frac{\bar{c}}{2}\sqrt{\Delta_{k_2}}\right)^2 + 2\bar{\delta}_1 
\end{equation}
where
\begin{equation*}
    \Delta_{k_1} =16\epsilon - 8 \left(\frac{(1+\epsilon)(k_M - k_m)}{\underbar{k}_1 k_m}\right) \times \left(\frac{(1+\epsilon)}{\underbar{k}_1k_m}2\bar{\delta}_1+ \frac{\alpha}{k_M}\right),
\end{equation*}
\begin{equation*}
    \bar{c} = \frac{\underbar{k}_1k_m}{(1+\epsilon)(k_M - k_m)}, \ \underbar{k}_1 = k_1 - \frac{\alpha}{k_m}
\end{equation*}
\end{corollary}

\subsection{Closing the loop}
Example revisited
\begin{equation*}
    \left(1 + \cos^2(q)\right)\ddot{q} + g\sin(q) + b\left(\dot{q} + \arctan (\dot{q})\right) = u
\end{equation*}

\subsection{Key takeaways}
Features
\begin{itemize}
    \item Deals well with matched uncertainties
    \item Works for uncertain input gain
    \item Rigorously proves control gain conditions for desired stability properties
\end{itemize}
Future work
\begin{itemize}
    \item Combine with adaptation law
\end{itemize}